\honest{Asch's Conformity Experiment}{Eliezer Yudkowsky}{2007}

In the 1950s Solomon Asch asked subjects which of three lines matched a standard line. Before the real subject answered, a row of confederates each gave the same wrong answer. ``Three-quarters of the subjects in Asch's experiment gave a `conforming' answer at least once. A third of the subjects conformed more than half the time.'' Afterwards most said they had not really believed their conforming answers, but some said they had. Asch found the strength of conformity ``a matter of concern.'' \nb{Both figures match Asch's 1956 table. Per answer, his subjects went along ``in 36.8 per cent of the selections,'' so nearly two answers in three were correct, and Asch wrote that the majority effect was not ``even the strongest force at work.'' In later American studies, conformity was lower than in Asch's.}

Was conforming irrational? Not obviously. Aumann's agreement theorem shows that honest Bayesians with common knowledge of their estimates must agree, so other people's beliefs are evidence. \nb{The theorem also requires that the two start from the same priors. The post leaves this condition out.} If three honest people who saw the same diagram all said C, I hope I would give the majority more than even odds, and say only that it looks to me like B, but I have no reason to think my judgment better than theirs. That is a much weaker claim than saying I see through the illusion that fools them. \nb{Some of Asch's independent subjects did this. They came to believe the majority was right, and kept reporting what they saw.}

Then I go through the details. Conformity rises up to three confederates and not beyond, though fifteen people should be stronger evidence than three. One dissenter, even one giving a different wrong answer, cuts conformity to ``5–10\% of subjects.'' If one against three means the three are probably right, one should equally trust six against two. But a nervous person made less alone by one ally fits the result easily. \nb{Asch's figures are shares of answers, not of subjects. A dissenter whose wrong answer was more extreme than the majority's cut errors to 9 per cent; one whose wrong answer was less extreme cut the majority's effect by only about a third.} Subjects deny that the dissenter helped them. Like the 90\% of drivers who think themselves above average, some may be right, but not all. People do not know the causes of their own conformity, which counts against calling the pattern rational. When the dissenter switches to the group, conformity returns. Being the first dissenter is a valuable and costly service, but ``you've got to keep it up.'' \nb{Asch put this down to the ``desertion.'' When the partner simply left the room, ``the partner's effect outlasted his presence.''}

In all-female groups about half the subjects conform more than half the time, against a third in all-male groups. If the average subject is rational, then women are too agreeable and men too disagreeable, and neither is rational. Members of an in-group, such as handicapped subjects among other handicapped people, conform more. Conformity is lower when the diagram is blatant, which is hard to explain if all subjects are strategically avoiding sticking out. It is lower when answers are private, which counts against Aumann.

\nb{The two readings are the distinction Deutsch and Gerard drew in 1955 between ``informational'' and ``normative'' social influence. The post does not mention it.}
